% Options for packages loaded elsewhere
\PassOptionsToPackage{unicode}{hyperref}
\PassOptionsToPackage{hyphens}{url}
\PassOptionsToPackage{dvipsnames,svgnames,x11names}{xcolor}
\documentclass[
  10pt,
]{article}
\usepackage{xcolor}
\usepackage[margin=1cm,top=1cm,bottom=1cm,left=1cm,right=1cm,includeheadfoot]{geometry}
\usepackage{amsmath,amssymb}
\setcounter{secnumdepth}{3}
\usepackage{iftex}
\ifPDFTeX
  \usepackage[T1]{fontenc}
  \usepackage[utf8]{inputenc}
  \usepackage{textcomp} % provide euro and other symbols
\else % if luatex or xetex
  \usepackage{unicode-math} % this also loads fontspec
  \defaultfontfeatures{Scale=MatchLowercase}
  \defaultfontfeatures[\rmfamily]{Ligatures=TeX,Scale=1}
\fi
\usepackage{lmodern}
\ifPDFTeX\else
  % xetex/luatex font selection
  \setmainfont[]{DejaVu Serif}
  \setmonofont[]{DejaVu Sans Mono}
\fi
% Use upquote if available, for straight quotes in verbatim environments
\IfFileExists{upquote.sty}{\usepackage{upquote}}{}
\IfFileExists{microtype.sty}{% use microtype if available
  \usepackage[]{microtype}
  \UseMicrotypeSet[protrusion]{basicmath} % disable protrusion for tt fonts
}{}
\usepackage{setspace}
\makeatletter
\@ifundefined{KOMAClassName}{% if non-KOMA class
  \IfFileExists{parskip.sty}{%
    \usepackage{parskip}
  }{% else
    \setlength{\parindent}{0pt}
    \setlength{\parskip}{6pt plus 2pt minus 1pt}}
}{% if KOMA class
  \KOMAoptions{parskip=half}}
\makeatother
\usepackage{listings}
\newcommand{\passthrough}[1]{#1}
\lstset{defaultdialect=[5.3]Lua}
\lstset{defaultdialect=[x86masm]Assembler}
\usepackage{graphicx}
\makeatletter
\newsavebox\pandoc@box
\newcommand*\pandocbounded[1]{% scales image to fit in text height/width
  \sbox\pandoc@box{#1}%
  \Gscale@div\@tempa{\textheight}{\dimexpr\ht\pandoc@box+\dp\pandoc@box\relax}%
  \Gscale@div\@tempb{\linewidth}{\wd\pandoc@box}%
  \ifdim\@tempb\p@<\@tempa\p@\let\@tempa\@tempb\fi% select the smaller of both
  \ifdim\@tempa\p@<\p@\scalebox{\@tempa}{\usebox\pandoc@box}%
  \else\usebox{\pandoc@box}%
  \fi%
}
% Set default figure placement to htbp
\def\fps@figure{htbp}
\makeatother
\setlength{\emergencystretch}{3em} % prevent overfull lines
\providecommand{\tightlist}{%
  \setlength{\itemsep}{0pt}\setlength{\parskip}{0pt}}
% Enable graphics inclusion and ensure figure numbering works
\usepackage{graphicx}
\renewcommand{\figurename}{Figure}

% Configure fonts for Unicode support with fallbacks
\usepackage{newunicodechar}
\newunicodechar{⁴}{\textsuperscript{4}}
\newunicodechar{₄}{\textsubscript{4}}

% Math notation macros
\newcommand{\norm}[1]{\lVert #1\rVert}
\newcommand{\abs}[1]{\lvert #1\rvert}

% Enhanced code block styling for better contrast and readability
\usepackage{fancyvrb}
\usepackage{xcolor}
\usepackage{listings}

% Define custom colors for code blocks
\definecolor{codebg}{RGB}{245, 245, 245}      % Light gray background
\definecolor{codeborder}{RGB}{200, 200, 200}  % Medium gray border
\definecolor{codefg}{RGB}{50, 50, 50}         % Dark gray text

% Configure Verbatim environment for inline code
\DefineVerbatimEnvironment{Verbatim}{Verbatim}{%
    fontsize=\small,
    frame=single,
    framerule=0.5pt,
    framesep=3pt,
    rulecolor=\color{codeborder},
    bgcolor=\color{codebg},
    fgcolor=\color{codefg}
}

% Configure code block styling
\DefineVerbatimEnvironment{Highlighting}{Verbatim}{%
    fontsize=\footnotesize,
    frame=single,
    framerule=0.5pt,
    framesep=5pt,
    rulecolor=\color{codeborder},
    bgcolor=\color{codebg},
    fgcolor=\color{codefg}
}

% Style inline code with \texttt
\renewcommand{\texttt}[1]{%
    \colorbox{codebg}{\color{codefg}\ttfamily #1}%
}

% Configure listings package for code blocks
\lstset{
    backgroundcolor=\color{codebg},
    basicstyle=\footnotesize\ttfamily\color{codefg},
    breakatwhitespace=false,
    breaklines=true,
    captionpos=b,
    commentstyle=\color{codefg},
    deletekeywords={...},
    escapeinside={\%*}{*)},
    extendedchars=true,
    frame=single,
    framerule=0.5pt,
    framesep=5pt,
    keepspaces=true,
    keywordstyle=\color{codefg},
    language=Python,
    morekeywords={*,...},
    numbers=left,
    numbersep=5pt,
    numberstyle=\tiny\color{codefg},
    rulecolor=\color{codeborder},
    showspaces=false,
    showstringspaces=false,
    showtabs=false,
    stepnumber=1,
    stringstyle=\color{codefg},
    tabsize=2,
    title=\lstname
}

% Override any Pandoc default lstset configurations
\AtBeginDocument{
    \lstset{
        backgroundcolor=\color{codebg},
        basicstyle=\footnotesize\ttfamily\color{codefg},
        frame=single,
        framerule=0.5pt,
        framesep=5pt,
        rulecolor=\color{codeborder},
        numbers=left,
        numbersep=5pt,
        numberstyle=\tiny\color{codefg}
    }
}

% Configure hyperref colors consistently
\AtBeginDocument{
% Override pandoc's hidelinks setting with consistent options
\hypersetup{
    colorlinks=true,
    allcolors=red,
    linkcolor=red,
    urlcolor=red,
    citecolor=red,
    filecolor=red,
    menucolor=red,
    linktoc=all
}
}

% Simple page break support for document structure
\usepackage{bookmark}
\IfFileExists{xurl.sty}{\usepackage{xurl}}{} % add URL line breaks if available
\urlstyle{same}
% fallback for those not using the hyperref driver hyperxmp:
\makeatletter
\@ifundefined{xmpquote}{\newcommand{\xmpquote}[1]{#1}}{}
\makeatother
\hypersetup{
  colorlinks=true,
  linkcolor={red},
  filecolor={red},
  citecolor={red},
  urlcolor={red},
  pdfcreator={LaTeX via pandoc}}

\title{Statistics and Scaling Gallery}
\author{Daniel Ari Friedman\\ ORCID: 0000-0001-6232-9096\\ Email: daniel@activeinference.institute\\ DOI: 10.5281/zenodo.16887791}
\date{October 07, 2026}

\begin{document}
\maketitle

{
\hypersetup{linkcolor=black}
\setcounter{tocdepth}{3}
\tableofcontents
}
\setstretch{1.0}
\section{Statistics and Scaling
Gallery}\label{statistics-and-scaling-gallery}

\subsection{Overview}\label{overview}

This section is the figure surface of the benchmarks-and-statistics
layer: every rendering primitive used by section
\passthrough{\lstinline!17\_benchmarks\_statistics.md!} lives in
\passthrough{\lstinline!src/quadmath/viz/vis\_stats.py!}
(\passthrough{\lstinline!plot\_latency\_hist!},
\passthrough{\lstinline!plot\_scaling\_loglog!},
\passthrough{\lstinline!plot\_ci\_bars!},
\passthrough{\lstinline!plot\_ecdf!}). The primitives are input-agnostic
--- plain numpy arrays in, artists out; they import nothing from the
other \passthrough{\lstinline!src/!} modules and receive their
matplotlib axes explicitly, so panels compose inside caller-owned
figures and nothing renders at import time. The command-line entry is
\passthrough{\lstinline!quadmath/scripts/stats\_gallery.py!} --- a thin
orchestrator that sets a headless backend and a fixed seed
(\passthrough{\lstinline!seed=32!}), delegates to the gallery function
in \passthrough{\lstinline!src/quadmath/viz/vis\_stats.py!}
(thin-orchestrator contract of
\passthrough{\lstinline!quadmath/scripts/AGENTS.md!}), and prints each
written path on its own line: those stdout lines are the
\passthrough{\lstinline!make\_all\_figures!} manifest contract. The
gallery writes exactly four PNGs ---
\passthrough{\lstinline!stats\_gallery\_latency.png!},
\passthrough{\lstinline!stats\_gallery\_scaling.png!},
\passthrough{\lstinline!stats\_gallery\_ci.png!},
\passthrough{\lstinline!stats\_gallery\_ecdf.png!} --- and every random
draw comes from the fixed seed, so the whole set is deterministic.

\subsection{Latency histogram}\label{latency-histogram}

\passthrough{\lstinline!plot\_latency\_hist(times, ax=None, *, bins=20, title=...)!}
renders the distribution of per-call wall-clock durations measured by
\passthrough{\lstinline!src/quadmath/stats/benchmarks.py::time\_callable!}
(trials summary in
\passthrough{\lstinline!17\_benchmarks\_statistics.md!}) as a histogram
with \passthrough{\lstinline!bins!} bins on the given axes. The shape of
the distribution carries what the mean alone hides: a tight spike means
the harness saw a stable workload, while a heavy right tail is exactly
what the \passthrough{\lstinline!median\_s!} /
\passthrough{\lstinline!p95\_s!} fields of
\passthrough{\lstinline!BenchRow!} are there to expose. The gallery
panel histograms a deterministic synthetic stand-in for such a sample:
256 draws of a lognormal distribution (mean 0, \(\sigma = 0.6\), clipped
to \([0.05, 5.0]\) in seconds-scale units) taken from the single
fixed-seed \passthrough{\lstinline!numpy.random.default\_rng(32)!}
generator, so re-renders are byte-identical.

\begin{figure}
\centering
\pandocbounded{\includegraphics[keepaspectratio,alt={Latency distribution. Histogram (20 bins) of a synthetic 256-draw lognormal latency sample (mean 0, \textbackslash sigma = 0.6) clipped to {[}0.05, 5.0{]} in seconds-scale units, drawn from the fixed-seed numpy.random.default\_rng(32) generator of stats\_gallery.py; the dashed line marks the sample mean. Rendered by plot\_latency\_hist; reproduced by uv run python quadmath/scripts/stats\_gallery.py. The spread and right tail complement the mean\_s, median\_s, and p95\_s fields of BenchRow, whose per-call times time\_callable measures with time.perf\_counter after one discarded warmup call.}]{../output/figures/stats_gallery_latency.png}}
\caption{\textbf{Latency distribution.} Histogram (20 bins) of a
synthetic 256-draw lognormal latency sample (mean 0, \(\sigma = 0.6\))
clipped to {[}0.05, 5.0{]} in seconds-scale units, drawn from the
fixed-seed
\protect\passthrough{\lstinline!numpy.random.default\_rng(32)!}
generator of \protect\passthrough{\lstinline!stats\_gallery.py!}; the
dashed line marks the sample mean. Rendered by
\protect\passthrough{\lstinline!plot\_latency\_hist!}; reproduced by
\protect\passthrough{\lstinline!uv run python quadmath/scripts/stats\_gallery.py!}.
The spread and right tail complement the
\protect\passthrough{\lstinline!mean\_s!},
\protect\passthrough{\lstinline!median\_s!}, and
\protect\passthrough{\lstinline!p95\_s!} fields of
\protect\passthrough{\lstinline!BenchRow!}, whose per-call times
\protect\passthrough{\lstinline!time\_callable!} measures with
\protect\passthrough{\lstinline!time.perf\_counter!} after one discarded
warmup call.}
\end{figure}

\subsection{Scaling fit}\label{scaling-fit}

\passthrough{\lstinline!plot\_scaling\_loglog(sizes, times, ax=None, *, title=...)!}
plots measured workload sizes against durations on log-log axes and
overlays the least-squares power law of
\passthrough{\lstinline!src/quadmath/stats/statistics.py::scaling\_fit!}
--- the slope is the empirical complexity exponent \(\beta_1\) of
\eqref{eq:stat-scaling}. On log-log axes a power law is a straight line,
so the panel shows at a glance whether a lattice routine scales
linearly, quadratically, or worse, and how tightly the data follow the
law (\(r^2\)). The gallery panel fits a synthetic sweep with known
ground truth: durations \(t \approx 3\,n^{1.35}\,(1 + \varepsilon)\)
with \(\varepsilon \sim \mathcal{N}(0, 0.02^{2})\) relative noise over
sizes \(n \in \{8, 16, 32, 64, 128, 256\}\), drawn from the same
fixed-seed generator (times rounded to four decimals for byte
stability), so the fit should recover the planted exponent 1.35.

\begin{figure}
\centering
\pandocbounded{\includegraphics[keepaspectratio,alt={Log-log scaling fit. Workload sizes n \textbackslash in \textbackslash\{8, 16, 32, 64, 128, 256\textbackslash\} (dimensionless units) against synthetic durations t \textbackslash approx 3\textbackslash,n\^{}\{1.35\}\textbackslash,(1 + \textbackslash varepsilon) with \textbackslash varepsilon \textbackslash sim \textbackslash mathcal\{N\}(0, 0.02\^{}\{2\}) relative noise, times rounded to four decimals, all drawn from the fixed-seed numpy.random.default\_rng(32) generator; the least-squares power law from scaling\_fit is overlaid by plot\_scaling\_loglog, and its fitted slope --- annotated in the legend --- should recover the planted exponent 1.35, the empirical complexity exponent \textbackslash beta\_1 of .}]{../output/figures/stats_gallery_scaling.png}}
\caption{\textbf{Log-log scaling fit.} Workload sizes
\(n \in \{8, 16, 32, 64, 128, 256\}\) (dimensionless units) against
synthetic durations \(t \approx 3\,n^{1.35}\,(1 + \varepsilon)\) with
\(\varepsilon \sim \mathcal{N}(0, 0.02^{2})\) relative noise, times
rounded to four decimals, all drawn from the fixed-seed
\protect\passthrough{\lstinline!numpy.random.default\_rng(32)!}
generator; the least-squares power law from
\protect\passthrough{\lstinline!scaling\_fit!} is overlaid by
\protect\passthrough{\lstinline!plot\_scaling\_loglog!}, and its fitted
slope --- annotated in the legend --- should recover the planted
exponent 1.35, the empirical complexity exponent \(\beta_1\) of
\eqref{eq:stat-scaling}.}
\end{figure}

\subsection{Confidence intervals}\label{confidence-intervals}

\passthrough{\lstinline!plot\_ci\_bars(labels, means, lows, highs, ax=None, *, title=...)!}
renders paired estimates as a bar-and-whisker chart: one bar per label
at its point estimate, with error bars spanning the low and high ends of
a percentile bootstrap interval
\passthrough{\lstinline!src/quadmath/stats/statistics.py::bootstrap\_ci!}
(\eqref{eq:stat-bootstrap}). Overlapping intervals visually encode the
same information a permutation test quantifies
(\eqref{eq:stat-permutation}): two conditions whose intervals do not
overlap are the ones the pooled test flags. The gallery panel shows
three benchmark estimates (\passthrough{\lstinline!to\_xyz!},
\passthrough{\lstinline!quadray\_from\_xyz!},
\passthrough{\lstinline!shell\_enum!}) with symmetric intervals whose
half-widths are seeded uniform draws on \([0.05, 0.25]\) from the same
fixed-seed \passthrough{\lstinline!numpy.random.default\_rng(32)!}
generator: the primitive renders whatever \passthrough{\lstinline!lows!}
/ \passthrough{\lstinline!highs!} it is given, and in a measured
analysis those bounds are the percentile endpoints
\passthrough{\lstinline!bootstrap\_ci!} returns
(\eqref{eq:stat-bootstrap}).

\begin{figure}
\centering
\pandocbounded{\includegraphics[keepaspectratio,alt={Point estimates with confidence intervals by condition. Three benchmark estimates (to\_xyz, quadray\_from\_xyz, shell\_enum) in synthetic estimate units, rendered by plot\_ci\_bars as points with symmetric error bars whose half-widths are uniform draws on {[}0.05, 0.25{]} from the fixed-seed numpy.random.default\_rng(32) generator; in a measured analysis the bar ends would be the 95\% percentile bootstrap endpoints from bootstrap\_ci (2000 resamples, ), whose non-overlap is the visual cue the permutation test of  quantifies.}]{../output/figures/stats_gallery_ci.png}}
\caption{\textbf{Point estimates with confidence intervals by
condition.} Three benchmark estimates
(\protect\passthrough{\lstinline!to\_xyz!},
\protect\passthrough{\lstinline!quadray\_from\_xyz!},
\protect\passthrough{\lstinline!shell\_enum!}) in synthetic estimate
units, rendered by \protect\passthrough{\lstinline!plot\_ci\_bars!} as
points with symmetric error bars whose half-widths are uniform draws on
{[}0.05, 0.25{]} from the fixed-seed
\protect\passthrough{\lstinline!numpy.random.default\_rng(32)!}
generator; in a measured analysis the bar ends would be the 95\%
percentile bootstrap endpoints from
\protect\passthrough{\lstinline!bootstrap\_ci!} (2000 resamples,
\eqref{eq:stat-bootstrap}), whose non-overlap is the visual cue the
permutation test of \eqref{eq:stat-permutation} quantifies.}
\end{figure}

\subsection{Empirical CDF}\label{empirical-cdf}

\passthrough{\lstinline!plot\_ecdf(values, ax=None, *, title=...)!}
renders the empirical distribution function

\begin{equation}
\label{eq:vis-ecdf}
\hat{F}(t) \;=\; \frac{1}{n}\sum_{i=1}^{n}
\mathbf{1}\{t_i \le t\},
\end{equation}

a monotone step function that jumps by \(1/n\) at every observed value
--- the exact distribution of the sample, with no binning choices at
all. For timing data the ECDF answers directly what a percentile summary
rounds off: the curve passes through the empirical median at height
\(0.5\) and through the 95th percentile at height \(0.95\), so
\passthrough{\lstinline!median\_s!} and \passthrough{\lstinline!p95\_s!}
are readable off the plot. The gallery panel draws the ECDF of the same
256-draw clipped lognormal latency sample as the histogram panel above,
so the two panels read as one distribution in two renderings.

\begin{figure}
\centering
\pandocbounded{\includegraphics[keepaspectratio,alt={Empirical cumulative distribution of the latency sample. Exact step-function ECDF \textbackslash hat\{F\}(t) of the same 256-draw clipped lognormal sample as the histogram panel above (fixed seed 32, latency in seconds-scale units), rendered by plot\_ecdf; the empirical median (height 0.5) and 95th percentile (height 0.95) are read directly off the curve.}]{../output/figures/stats_gallery_ecdf.png}}
\caption{\textbf{Empirical cumulative distribution of the latency
sample.} Exact step-function ECDF \(\hat{F}(t)\) of the same 256-draw
clipped lognormal sample as the histogram panel above (fixed seed 32,
latency in seconds-scale units), rendered by
\protect\passthrough{\lstinline!plot\_ecdf!}; the empirical median
(height 0.5) and 95th percentile (height 0.95) are read directly off the
curve.}
\end{figure}

\subsection{Reproducibility and test
contract}\label{reproducibility-and-test-contract}

\begin{itemize}
\tightlist
\item
  The gallery function in
  \passthrough{\lstinline!src/quadmath/viz/vis\_stats.py!} writes
  exactly the four files listed above, in that order, into the output
  directory (creating it when missing) and returns their paths. Every
  panel is fully seeded --- the script fixes
  \passthrough{\lstinline!seed=32!} --- and the PNGs carry no timestamp
  metadata, so re-renders are byte-identical.
\item
  Tests: \passthrough{\lstinline!tests/test\_vis\_stats.py!} pins the
  contract with deterministic assertions, no mocks, and no pixel diffs
  --- artist placement on caller-provided axes, input-agnostic behavior
  over plain arrays, and the invariants of each primitive (monotone
  non-decreasing ECDF steps, CI whisker endpoints matching the
  lows/highs inputs). Consistent with
  \passthrough{\lstinline!17\_benchmarks\_statistics.md!}, no test
  asserts wall-clock timing values.
\item
  Figure regeneration:
  \passthrough{\lstinline!uv run python quadmath/scripts/stats\_gallery.py!}
  (the script sets \passthrough{\lstinline!MPLBACKEND=Agg!} itself);
  stdout is exactly the four written paths, the same manifest contract
  as \passthrough{\lstinline!quadmath/scripts/lattice\_gallery.py!}.
\item
  Markdown:
  \passthrough{\lstinline!uv run python quadmath/scripts/validate\_markdown.py!}.
\end{itemize}

\subsection{Cross-references}\label{cross-references}

\begin{itemize}
\tightlist
\item
  The measured quantities behind every panel:
  \passthrough{\lstinline!17\_benchmarks\_statistics.md!}
  (\passthrough{\lstinline!src/quadmath/stats/benchmarks.py!},
  \passthrough{\lstinline!src/quadmath/stats/statistics.py!}).
\item
  The sibling gallery of the lattice layer, same thin-orchestrator and
  determinism contract:
  \passthrough{\lstinline!16\_lattice\_gallery.md!}
  (\passthrough{\lstinline!src/quadmath/viz/vis\_lattice.py!}).
\end{itemize}

\end{document}
